\section{Usability, Fault Tolerance, and Operational Scope}
\label{sec:usability_and_overhead}

\subsection{Analytical Cognitive Latency Envelope}
Cognitive authentication architectures must remain within acceptable human completion bounds while preserving narrow cryptographic expiration windows~\cite{kepkowski2023challenges, binbeshr2025challenge}. Prior to human-subject testing, we establish a literature-derived additive latency envelope adapted from grid-based challenge-response studies~\cite{binbeshr2025challenge, duzgun2022shoulder}:
\begin{equation}
T_{total} = T_{perceptual} + T_{cognitive} + T_{motor}
\label{eq:latency}
\end{equation}
Perceptual search $T_{perceptual}$ spans $0.35\text{ s}$ to $0.60\text{ s}$ to locate $N_{ch}$ inside the 4-cell quadrant $q_u$. Cognitive traversal $T_{cognitive}$ estimates mental execution across four toroidal offsets at $450\text{ ms}$ to $725\text{ ms}$ per step ($1.80\text{ s}$ to $2.90\text{ s}$)~\cite{binbeshr2025challenge}. Physical input $T_{motor}$ for a 4-character string averages $1.10\text{ s}$ to $1.50\text{ s}$~\cite{incel2021dakota, khan2024mouse}. Aggregating these components yields an analytical mean target of $3.25\text{ s}$ to $5.00\text{ s}$, comfortably inside the 15-second TTL. Because toroidal wrap-around ($\bmod 4$) may impose higher mental load during initial onboarding than linear grid lookups, empirical validation across trained and novice cohorts is prioritized in Section~\ref{sec:experimental_design_and_thesis}.

\subsection{Error Rates, Retention, and Hardware-Free Recovery}
Empirical evaluations of shoulder-surfing-resistant challenge-response and graphical PIN schemes report initial human entry error rates between $4\%$ and $8\%$, alongside recall latency growth after multi-week inactivity~\cite{binbeshr2025challenge, duzgun2022shoulder}. To balance fault tolerance with intersection safety, BC-MFA enforces a layered verification policy. Upon a single algebraic mismatch ($R_t \ne R_{expected}$), the IdP issues one immediate retry on a fresh matrix $M_{t+1}$ without incrementing the untrusted intersection counter, provided behavioral telemetry confirms human liveness (failed submissions leak no valid $(M_t, R_t)$ pair to scrapers). If hesitation exceeds $\tau_{cog} > 6.0\text{ s}$, the UI can display a user-enrolled semantic mnemonic cue that aids path recall without revealing grid coordinates.

If response $R_t$ is valid but the behavioral score falls within a marginal uncertainty band ($\theta_{bio} - \epsilon \le \text{Score}(\mathcal{B}) < \theta_{bio}$), the IdP requests one confirmatory matrix rather than locking the session. After two consecutive cognitive failures, rather than falling back to vulnerable SMS OTP or mandating a corporate managed phone (which would contradict hardware independence and reopen recovery flaws documented by Amft et al.~\cite{amft2023weve}), the IdP enforces exponential cooldown throttling and routes account recovery to policy-selected strong channels: trusted-device re-authentication, pre-printed single-use cryptographic recovery codes, or electronic identity (eID) verification following Schwarz et al.~\cite{schwarz2022feido}.

\subsection{Enrollment Bias, Accessibility, and Application Scope}
Allowing users to choose their own offset paths induces predictable spatial patterns (e.g., linear horizontal sweeps) that shrink effective secret entropy well below $\Omega = 32{,}760$. BC-MFA therefore employs server-assisted uniform random secret generation paired with interactive visual mnemonic onboarding on a trusted device. Furthermore, because visual-spatial modular arithmetic under a 15-second window creates barriers for users with visual, motor, or cognitive disabilities, enterprise deployments must treat BC-MFA as an optional step-up authenticator paired with WCAG-compliant alternatives (such as hardware passkeys or audio-tactile schemes) rather than a universal mandatory replacement~\cite{mahfouz2025passkeys}. Its primary operational scope is targeted, hardware-free Step-Up defense on unmanaged BYOD endpoints, shared university or hospital terminals, and public kiosks where resident user-space infostealers threaten routine password or OTP entry.
